% Lösungen Abstände (zusammenbau v0.4, Heft)
\einheitenkopf{Abstände}

\erg{65}{a) Punkt–Punkt: Betrag des Verbindungsvektors \quad b) $\overrightarrow{AB} = (2 | 2 | 1)$, $d = \sqrt{4 + 4 + 1} = 3$ \quad c) $|\overrightarrow{AB}| = |(600 | 1\,200 | -400)| = 1\,400$ m, $v = 1\,400 : 40 = 35$ m/s $= 126$ km/h}
\erg{66}{a) $|\overrightarrow{PQ}| = |(60 | 80 | 0)| = 100$ m, $v = 100 : 4 = 25$ m/s $= 90$ km/h \quad b) $M_1(2 | 1 | 4)$, $M_2(4 | 7 | 1)$, $|\overrightarrow{M_1M_2}| = \sqrt{4 + 36 + 9} = 7$ m; Seil $7 \cdot 1{,}3 = 9{,}1$ m \quad c) Entfernteste Ecke ist C: $|\overrightarrow{PC}| = \sqrt{36 + 25 + 4} \approx 8{,}06$ m $< 9$ m – ja, jede Stelle ist erreichbar \quad d) Katheten gleich: $|\overrightarrow{AB}| = 3$, $|\overrightarrow{AC}| = \sqrt{5}\,t$, also $t = \tfrac{3}{\sqrt{5}} \approx 1{,}34$ \quad e) $|\overrightarrow{PQ}| = \sqrt{64 + (y-3)^2 + 225} \ge \sqrt{289} = 17$, weil $(y-3)^2 \ge 0$ für jedes y \quad f) $|\overrightarrow{RS}| = \sqrt{25 + 144 + (z-2)^2} \ge \sqrt{169} = 13$, weil $(z-2)^2 \ge 0$ für jedes z}
\erg{67}{a) Abstand berechnen \quad b) $|\vec n| = 3$, $d = \tfrac{|8 - 0 + 6 - 5|}{3} = 3$ \quad c) $\vec x = (-1 | 0 | 5) + t \cdot (3 | 0 | -4)$}
\erg{68}{a) P hat von E den halben Abstand 7; $|\vec n| = 5$: $P = (2 | 1 | 0) + \tfrac{7}{5} \cdot (0 | 3 | 4) = (2 | 5{,}2 | 5{,}6)$ (oder (2 | −3,2 | −5,6))}
\erg{69}{a) Lotbedingung (Skalarprodukt null) \quad b) $F_t(2 + 2t | 1 - t | 2t)$, $\overrightarrow{PF_t} \circ (2 | -1 | 2) = 9t - 9 = 0$, $t = 1$: F(4 | 0 | 2); Abstand $|\overrightarrow{PF}| = |(0 | -2 | -1)| \approx 2{,}24$ \quad c) X ist der Lotfußpunkt von P auf die Gerade durch M mit der Richtung $\vec u$; $|\overrightarrow{PX}|$ ist der Abstand von P zu dieser Geraden.}
\erg{70}{a) Q ist der Lotfußpunkt des Lots von E auf die Gerade durch A und B: I legt Q auf diese Gerade, II macht $\overrightarrow{EQ}$ senkrecht zu ihr. \quad b) M als Mittelpunkt von AB berechnen; R als allgemeinen Punkt $R_t$ der Geraden CD ansetzen; das Skalarprodukt von $\overrightarrow{MR_t}$ und $\overrightarrow{CD}$ null setzen und nach t lösen; t einsetzen – R ist der Lotfußpunkt, nicht der Mittelpunkt von CD. \quad c) I: Gerade durch Q in der xy-Ebene, senkrecht zur Kante AB (Richtung (1 | −2 | 0)); II: Schnitt mit L, in der xy-Ebene also mit der Geraden AB – R ist der Lotfußpunkt von Q auf AB; III: der Abstand von Q zur Geraden AB, $|\overrightarrow{QR}| \approx 2{,}24$ \quad d) P ist ein Punkt der Kante BE; $\overrightarrow{PA} \circ \overrightarrow{PB} = 0$ fordert, dass PA senkrecht auf der Kante steht – so ist der Weg von A zur Kante am kürzesten; weil A und C symmetrisch zur Ebene durch B, D, E liegen, ist der Weg von P nach C gleich lang, daher der Faktor 2: kürzeste Linie $\approx 5{,}96$ \quad e) Skizze: rechtwinkliges Dreieck mit der Kathete 0,3 (Abstand zu k) und der Hypotenuse QS: $|QS| = \tfrac{0{,}3}{\mathrm{sin}\,30^\circ} = 0{,}6$; $S = Q - 0{,}6 \cdot \tfrac{1}{5}(3 | 4 | 0) = (5{,}64 | 7{,}52 | 0)$}
\erg{71}{a) senkrecht: die Strecke ist der Abstand \quad b) größer: $\overrightarrow{PQ} = (1 | 2 | 3)$ ist kein Vielfaches von (2 | −1 | 2), PQ steht schräg; das Lot ist kürzer (Kontrolle: $|\overrightarrow{PQ}| \approx 3{,}74$, Abstand der Scheiben $= 2$) \quad c) Das Lot von F auf t trifft t in L(2,5 | 5,5) oberhalb von G, außerhalb der Strecke GJ; nur dort wird der Abstand angenommen, alle Punkte der Rampe sind weiter von F entfernt.}

71c)

\begin{ksys}[xmin=0,xmax=8,ymin=0,ymax=8,klein]
\gerade{-1}{8}{t}
\punkt{4}{4}{G}
\punkt{7}{1}{J}
\punkt{0}{3}{F}
\gerade{1}{3}{}
\punkt{2.5}{5.5}{L}
\end{ksys}

\erg{72}{a) größer: $\overrightarrow{BE} = (0 | 0 | -3)$ ist kein Vielfaches des Normalenvektors (1 | 2 | 2), BE steht nicht senkrecht auf den Ebenen, und das Lot ist die kürzeste Verbindung (Kontrolle: $|\overrightarrow{BE}| = 3$, Ebenenabstand $= 2$) \quad b) Der Abstand ist die Länge des Lots von F auf die Ebene; FG ist schräg, also Hypotenuse eines rechtwinkligen Dreiecks mit dem Lot als Kathete – das Lot ist kürzer als FG. \quad c) Seilstrecke, horizontaler Abstand und Höhenunterschied bilden ein rechtwinkliges Dreieck; die Seilstrecke ist die Hypotenuse, der horizontale Abstand eine Kathete, also kürzer. \quad d) Rechter Winkel bei P, also ist der Mittelpunkt M(5 | 3 | 0) der Hypotenuse QR der Umkreismittelpunkt (Thaleskreis); alle Punkte (5 | 3 | h) auf dem Lot durch M sind von P, Q, R gleich weit entfernt, etwa M und (5 | 3 | 5) \quad e) F und G liegen 3 unter der Spitze, $Q_1$ und $Q_2$ in der xy-Ebene 9 darunter; FG und $Q_1Q_2$ sind beide waagerecht, also parallel; Strahlensatz mit Zentrum S: $|Q_1Q_2| : |FG| = 9 : 3 = 3$ \quad f) F und G liegen 2 unter der Spitze, die xy-Ebene 8 darunter; FG und $Q_1Q_2$ sind parallel (beide waagerecht); Strahlensatz: $|Q_1Q_2| : |FG| = 8 : 2 = 4$}
